\documentclass[10pt,a4paper,onecolumn]{article}
\usepackage{marginnote}
\usepackage{graphicx}
\usepackage{xcolor}
\usepackage{authblk,etoolbox}
\usepackage{titlesec}
\usepackage{calc}
\usepackage{tikz}
\usepackage{hyperref}
\hypersetup{colorlinks,breaklinks=true,
            urlcolor=[rgb]{0.0, 0.5, 1.0},
            linkcolor=[rgb]{0.0, 0.5, 1.0}}
\usepackage{caption}
\usepackage{tcolorbox}
\usepackage{amssymb,amsmath}
\usepackage{seqsplit}
\usepackage{xstring}

\usepackage{float}
\let\origfigure\figure
\let\endorigfigure\endfigure
\renewenvironment{figure}[1][2] {
    \expandafter\origfigure\expandafter[H]
} {
    \endorigfigure
}

\usepackage[
  backend=biber,
%  style=alphabetic,
%  citestyle=numeric
]{biblatex}
\bibliography{paper.bib}

% --- Splitting \texttt --------------------------------------------------

\let\textttOrig=\texttt
\def\texttt#1{\expandafter\textttOrig{\seqsplit{#1}}}
\renewcommand{\seqinsert}{\ifmmode
  \allowbreak
  \else\penalty6000\hspace{0pt plus 0.02em}\fi}


% --- Pandoc does not distinguish between links like [foo](bar) and
% --- [foo](foo) -- a simplistic Markdown model.  However, this is
% --- wrong:  in links like [foo](foo) the text is the url, and must
% --- be split correspondingly.
% --- Here we detect links \href{foo}{foo}, and also links starting
% --- with https://doi.org, and use path-like splitting (but not
% --- escaping!) with these links.
% --- Another vile thing pandoc does is the different escaping of
% --- foo and bar.  This may confound our detection.
% --- This problem we do not try to solve at present, with the exception
% --- of doi-like urls, which we detect correctly.


\makeatletter
\let\href@Orig=\href
\def\href@Urllike#1#2{\href@Orig{#1}{\begingroup
    \def\Url@String{#2}\Url@FormatString
    \endgroup}}
\def\href@Notdoi#1#2{\def\tempa{#1}\def\tempb{#2}%
  \ifx\tempa\tempb\relax\href@Urllike{#1}{#2}\else
  \href@Orig{#1}{#2}\fi}
\def\href#1#2{%
  \IfBeginWith{#1}{https://doi.org}%
  {\href@Urllike{#1}{#2}}{\href@Notdoi{#1}{#2}}}
\makeatother

\newlength{\cslhangindent}
\setlength{\cslhangindent}{1.5em}
\newlength{\csllabelwidth}
\setlength{\csllabelwidth}{3em}
\newenvironment{CSLReferences}[3] % #1 hanging-ident, #2 entry spacing
 {% don't indent paragraphs
  \setlength{\parindent}{0pt}
  % turn on hanging indent if param 1 is 1
  \ifodd #1 \everypar{\setlength{\hangindent}{\cslhangindent}}\ignorespaces\fi
  % set entry spacing
  \ifnum #2 > 0
  \setlength{\parskip}{#2\baselineskip}
  \fi
 }%
 {}
\usepackage{calc}
\newcommand{\CSLBlock}[1]{#1\hfill\break}
\newcommand{\CSLLeftMargin}[1]{\parbox[t]{\csllabelwidth}{#1}}
\newcommand{\CSLRightInline}[1]{\parbox[t]{\linewidth - \csllabelwidth}{#1}}
\newcommand{\CSLIndent}[1]{\hspace{\cslhangindent}#1}

% --- Page layout -------------------------------------------------------------
\usepackage[top=3.5cm, bottom=3cm, right=1.5cm, left=1.0cm,
            headheight=2.2cm, reversemp, includemp, marginparwidth=4.5cm]{geometry}

% --- Default font ------------------------------------------------------------
\renewcommand\familydefault{\sfdefault}

% --- Style -------------------------------------------------------------------
\renewcommand{\bibfont}{\small \sffamily}
\renewcommand{\captionfont}{\small\sffamily}
\renewcommand{\captionlabelfont}{\bfseries}

% --- Section/SubSection/SubSubSection ----------------------------------------
\titleformat{\section}
  {\normalfont\sffamily\Large\bfseries}
  {}{0pt}{}
\titleformat{\subsection}
  {\normalfont\sffamily\large\bfseries}
  {}{0pt}{}
\titleformat{\subsubsection}
  {\normalfont\sffamily\bfseries}
  {}{0pt}{}
\titleformat*{\paragraph}
  {\sffamily\normalsize}


% --- Header / Footer ---------------------------------------------------------
\usepackage{fancyhdr}
\pagestyle{fancy}
\fancyhf{}
%\renewcommand{\headrulewidth}{0.50pt}
\renewcommand{\headrulewidth}{0pt}
\fancyhead[L]{\hspace{-0.75cm}\includegraphics[width=5.5cm]{./logo.png}}
\fancyhead[C]{}
\fancyhead[R]{}
\renewcommand{\footrulewidth}{0.25pt}

\fancyfoot[L]{\parbox[t]{0.98\headwidth}{\footnotesize{\sffamily ¿citation\_author?, (2025). The
Population Synthesis Toolkit (PST) Python
Library. \textit{Journal of Open Source Software}, ¿VOL?(¿ISSUE?), ¿PAGE?. \url{https://doi.org/DOI unavailable}}}}


\fancyfoot[R]{\sffamily \thepage}
\makeatletter
\let\ps@plain\ps@fancy
\fancyheadoffset[L]{4.5cm}
\fancyfootoffset[L]{4.5cm}

% --- Macros ---------

\definecolor{linky}{rgb}{0.0, 0.5, 1.0}

\newtcolorbox{repobox}
   {colback=red, colframe=red!75!black,
     boxrule=0.5pt, arc=2pt, left=6pt, right=6pt, top=3pt, bottom=3pt}

\newcommand{\ExternalLink}{%
   \tikz[x=1.2ex, y=1.2ex, baseline=-0.05ex]{%
       \begin{scope}[x=1ex, y=1ex]
           \clip (-0.1,-0.1)
               --++ (-0, 1.2)
               --++ (0.6, 0)
               --++ (0, -0.6)
               --++ (0.6, 0)
               --++ (0, -1);
           \path[draw,
               line width = 0.5,
               rounded corners=0.5]
               (0,0) rectangle (1,1);
       \end{scope}
       \path[draw, line width = 0.5] (0.5, 0.5)
           -- (1, 1);
       \path[draw, line width = 0.5] (0.6, 1)
           -- (1, 1) -- (1, 0.6);
       }
   }

% --- Title / Authors ---------------------------------------------------------
% patch \maketitle so that it doesn't center
\patchcmd{\@maketitle}{center}{flushleft}{}{}
\patchcmd{\@maketitle}{center}{flushleft}{}{}
% patch \maketitle so that the font size for the title is normal
\patchcmd{\@maketitle}{\LARGE}{\LARGE\sffamily}{}{}
% patch the patch by authblk so that the author block is flush left
\def\maketitle{{%
  \renewenvironment{tabular}[2][]
    {\begin{flushleft}}
    {\end{flushleft}}
  \AB@maketitle}}
\makeatletter
\renewcommand\AB@affilsepx{ \protect\Affilfont}
%\renewcommand\AB@affilnote[1]{{\bfseries #1}\hspace{2pt}}
\renewcommand\AB@affilnote[1]{{\bfseries #1}\hspace{3pt}}
\renewcommand{\affil}[2][]%
   {\newaffiltrue\let\AB@blk@and\AB@pand
      \if\relax#1\relax\def\AB@note{\AB@thenote}\else\def\AB@note{#1}%
        \setcounter{Maxaffil}{0}\fi
        \begingroup
        \let\href=\href@Orig
        \let\texttt=\textttOrig
        \let\protect\@unexpandable@protect
        \def\thanks{\protect\thanks}\def\footnote{\protect\footnote}%
        \@temptokena=\expandafter{\AB@authors}%
        {\def\\{\protect\\\protect\Affilfont}\xdef\AB@temp{#2}}%
         \xdef\AB@authors{\the\@temptokena\AB@las\AB@au@str
         \protect\\[\affilsep]\protect\Affilfont\AB@temp}%
         \gdef\AB@las{}\gdef\AB@au@str{}%
        {\def\\{, \ignorespaces}\xdef\AB@temp{#2}}%
        \@temptokena=\expandafter{\AB@affillist}%
        \xdef\AB@affillist{\the\@temptokena \AB@affilsep
          \AB@affilnote{\AB@note}\protect\Affilfont\AB@temp}%
      \endgroup
       \let\AB@affilsep\AB@affilsepx
}
\makeatother
\renewcommand\Authfont{\sffamily\bfseries}
\renewcommand\Affilfont{\sffamily\small\mdseries}
\setlength{\affilsep}{1em}


\ifnum 0\ifxetex 1\fi\ifluatex 1\fi=0 % if pdftex
  \usepackage[T1]{fontenc}
  \usepackage[utf8]{inputenc}

\else % if luatex or xelatex
  \ifxetex
    \usepackage{mathspec}
    \usepackage{fontspec}

  \else
    \usepackage{fontspec}
  \fi
  \defaultfontfeatures{Ligatures=TeX,Scale=MatchLowercase}

\fi
% use upquote if available, for straight quotes in verbatim environments
\IfFileExists{upquote.sty}{\usepackage{upquote}}{}
% use microtype if available
\IfFileExists{microtype.sty}{%
\usepackage{microtype}
\UseMicrotypeSet[protrusion]{basicmath} % disable protrusion for tt fonts
}{}

\usepackage{hyperref}
\hypersetup{unicode=true,
            pdftitle={The Population Synthesis Toolkit (PST) Python Library},
            pdfborder={0 0 0},
            breaklinks=true}
\urlstyle{same}  % don't use monospace font for urls

% --- We redefined \texttt, but in sections and captions we want the
% --- old definition
\let\addcontentslineOrig=\addcontentsline
\def\addcontentsline#1#2#3{\bgroup
  \let\texttt=\textttOrig\addcontentslineOrig{#1}{#2}{#3}\egroup}
\let\markbothOrig\markboth
\def\markboth#1#2{\bgroup
  \let\texttt=\textttOrig\markbothOrig{#1}{#2}\egroup}
\let\markrightOrig\markright
\def\markright#1{\bgroup
  \let\texttt=\textttOrig\markrightOrig{#1}\egroup}


\IfFileExists{parskip.sty}{%
\usepackage{parskip}
}{% else
\setlength{\parindent}{0pt}
\setlength{\parskip}{6pt plus 2pt minus 1pt}
}
\setlength{\emergencystretch}{3em}  % prevent overfull lines
\providecommand{\tightlist}{%
  \setlength{\itemsep}{0pt}\setlength{\parskip}{0pt}}
\setcounter{secnumdepth}{0}
% Redefines (sub)paragraphs to behave more like sections
\ifx\paragraph\undefined\else
\let\oldparagraph\paragraph
\renewcommand{\paragraph}[1]{\oldparagraph{#1}\mbox{}}
\fi
\ifx\subparagraph\undefined\else
\let\oldsubparagraph\subparagraph
\renewcommand{\subparagraph}[1]{\oldsubparagraph{#1}\mbox{}}
\fi

\title{The Population Synthesis Toolkit (PST) Python Library}

        \author[1]{Pablo Corcho-Caballero}
          \author[2, 3]{Yago Ascasibar}
          \author[4]{Daniel Jim\'enez-L\'opez}
    
      \affil[1]{Kapteyn Astronomical Institute, University of Groningen,
The Netherlands}
      \affil[2]{Department of Theoretical Physics, Universidad Aut\'onoma
de Madrid (UAM), Spain}
      \affil[3]{Centro de Investigaci\'on Avanzada en F\'isica Fundamental
(CIAFF-UAM), Spain}
      \affil[4]{Observatorio Astron\'omico Nacional (OAN), Spain}
  \date{\vspace{-7ex}}

\begin{document}
\maketitle

\marginpar{

  \begin{flushleft}
  %\hrule
  \sffamily\small

  {\bfseries DOI:} \href{https://doi.org/DOI unavailable}{\color{linky}{DOI unavailable}}

  \vspace{2mm}

  {\bfseries Software}
  \begin{itemize}
    \setlength\itemsep{0em}
    \item \href{N/A}{\color{linky}{Review}} \ExternalLink
    \item \href{NO_REPOSITORY}{\color{linky}{Repository}} \ExternalLink
    \item \href{DOI unavailable}{\color{linky}{Archive}} \ExternalLink
  \end{itemize}

  \vspace{2mm}

  \par\noindent\hrulefill\par

  \vspace{2mm}

  {\bfseries Editor:} \href{https://example.com}{Pending
Editor} \ExternalLink \\
  \vspace{1mm}
    {\bfseries Reviewers:}
  \begin{itemize}
  \setlength\itemsep{0em}
    \item \href{https://github.com/Pending Reviewers}{@Pending
Reviewers}
    \end{itemize}
    \vspace{2mm}

  {\bfseries Submitted:} N/A\\
  {\bfseries Published:} N/A

  \vspace{2mm}
  {\bfseries License}\\
  Authors of papers retain copyright and release the work under a Creative Commons Attribution 4.0 International License (\href{http://creativecommons.org/licenses/by/4.0/}{\color{linky}{CC BY 4.0}}).

  
  \end{flushleft}
}

\hypertarget{summary}{%
\section{Summary}\label{summary}}

Stellar population synthesis is a crucial methodology in astrophysics,
enabling the interpretation of the integrated light of galaxies and
stellar clusters. By combining empirical and/or theoretical libraries of
the spectral energy distributions emitted by simple stellar populations
(SSPs) with star formation history (SFH) and chemical evolution models,
population synthesis can help estimate essential properties of galaxies,
such as total stellar mass, star formation rate, mass-weighted age,
metallicity, and so on.

PST is a Python library that offers a comprehensive and flexible
framework for stellar population synthesis. Its main goal is to easily
and efficiently compute composite spectra using different galaxy
evolution models and SSP libraries. It also incorporates additional
effects such as cosmic redshift, dust extinction and attenuation, and
computes several observable quantities derived from the spectra,
including broadband photometric fluxes and equivalent widths.

\hypertarget{state-of-the-field}{%
\section{State of the field}\label{state-of-the-field}}

A number of software packages have been developed to support stellar
population synthesis and modeling of galaxy spectral energy
distributions. Tools such as
\href{https://gitlab.surrey.ac.uk/ri0005/binary_c-python}{binary\_c-python}
(Hendriks \& Izzard, 2023) and
\href{https://github.com/astropy/SPISEA}{SPISEA} (Hosek et al., 2020)
are designed primarily for generating and analyzing simple stellar
populations, often with a focus on individual stars, binaries, or star
clusters. Meanwhile, libraries such as
\href{https://github.com/dfm/python-fsps}{python-FSPS} (Johnson et al.,
2024), a Python interface to the Flexible Stellar Population Synthesis
(FSPS) code (Conroy et al., 2009; Conroy \& Gunn, 2010), and the more
recent \href{https://github.com/ArgonneCPAC/dsps}{DSPS} (Hearin et al.,
2023), implemented using JAX for efficient gradient computation and
forward modeling, provide extensive modeling capabilities, although they
are sometimes limited to a specific set of SSP models or isochrones.

Other packages put a stronger emphasis on fitting observed data to
derive galaxy properties. Examples of these include Bayesian frameworks
such as \href{https://cigale.lam.fr/}{CIGALE} (Boquien et al., 2019),
\href{https://github.com/asgr/ProSpect}{ProSpect} (Robotham et al.,
2020) and \href{https://prospect.readthedocs.io/en/v1.0.0/}{Prospector}
(Johnson et al., 2021), which infer star formation histories and other
physical parameters using spectro-photometric data. Alternative
frequentist tools such as \href{https://pypi.org/project/ppxf/}{PpXF}
(Cappellari \& Emsellem, 2004),
\href{http://www.starlight.ufsc.br/}{Starlight} (Cid Fernandes et al.,
2005), or \href{https://gitlab.com/pipe3d/pyPipe3D}{Pipe3D} (S\'anchez et
al., 2016), are commonly used to extract stellar kinematics and stellar
population parameters from observed galaxy spectra, often in the context
of integral field spectroscopy.

\hypertarget{statement-of-need}{%
\section{Statement of need}\label{statement-of-need}}

The user-friendly modular framework of PST is designed to address the
following challenges:

\begin{itemize}
\tightlist
\item
  Handle a broad variety of SSP libraries that are publicly available in
  heterogeneous native formats.
\item
  Model arbitrarily complex galaxy star formation and chemical evolution
  histories.
\item
  Enable the simultaneous and self-consistent analysis of photometric
  and spectroscopic data from different instruments.
\end{itemize}

PST is designed for astronomy researchers, especially those working in
extragalactic astrophysics and stellar population synthesis, who require
a flexible and extensible Python-based toolkit for modeling galaxy
properties. PST is suited for users with intermediate to advanced Python
expertise and familiarity with common data formats and concepts in
astronomical spectroscopy and photometry.

The primary use cases are data analysis, synthetic model construction,
and pipeline integration for studies involving stellar population
synthesis (see the examples below). PST is particularly useful in
workflows that combine observational data with theoretical models within
a Bayesian or forward-modeling framework.

PST is currently a dependency of
\href{https://github.com/pykoala}{PyKOALA} (Corcho-Caballero et al.,
2025), another open-source Python package focused on reducing optical
integral-field spectroscopic observations. There, PST is mainly used to
derive broadband photometry. PST is also at the core of the Bayesian
Estimator for Stellar Population Analysis
(\href{https://https://besta.readthedocs.io/}{BESTA}, see also Euclid
Collaboration: Corcho-Caballero et al., 2025), where it is coupled with
the \href{https://cosmosis.readthedocs.io/en/latest/}{CosmoSIS} (Zuntz
et al., 2015) Monte Carlo sampling framework to infer the physical
properties of galaxies from the observed colors and spectra.

\hypertarget{features-and-functionality}{%
\section{Features and functionality}\label{features-and-functionality}}

PST design is built around three main components. First, the SSP module
enables the consistent use and manipulation of different SSP libraries.
This allows for the seamless ingestion of models and data from various
literature sources. The current version includes interfaces to a range
of SSP models, including PopStar (Moll\'a et al., 2009), Bruzual and
Charlot (BC03, Bruzual \& Charlot, 2003), E-MILES (Vazdekis et al.,
2016), and the X-Shooter Spectral Library (XSL, Verro et al., 2022) SSP
models.

For any SSP model integrated into PST, the library provides tools for
interpolating across stellar ages, metallicities, and wavelengths. Users
can easily compute key SSP quantities, such as the stellar mass-to-light
ratio in a given band, colors, and line indices.

Second, the \texttt{ChemicalEvolutionModel} classes represent the star
formation and chemical enrichment histories required to produce
composite spectral energy distributions and additional derived
quantities. These classes implement several widely-used analytic
prescriptions for modeling star formation histories (SFHs), such as
exponentially-declining or log-normal models. They also implement
complex SFH representations, such as table-based SFHs and particle-like
data models. These models are particularly suitable for post-processing
results from cosmological hydrodynamic simulations.

Third, PST features a dedicated \texttt{observables} module that can
predict additional quantities from spectra, such as broadband
photometric fluxes, colors, and equivalent widths, which are useful for
estimating the strength of absorption or emission lines. PST includes
automatic integration with the photometric filters provided by the
\href{http://svo2.cab.inta-csic.es/theory/fps/}{Spanish Virtual
Observatory Filter Profile Service} (Rodrigo et al., 2024, 2012; Rodrigo
\& Solano, 2020) for synthetic photometry calculations, as well as
popular line indices such as the Lick system (Worthey et al., 1994).

\hypertarget{documentation-and-tutorials}{%
\section{Documentation and
tutorials}\label{documentation-and-tutorials}}

To make PST easier to use, we provide a set of comprehensive tutorials
in the form of Jupyter notebooks. These tutorials cover the following
topics:

\begin{itemize}
\tightlist
\item
  Interacting with SSP models and exploring their fundamental
  properties.
\item
  Producing composite spectra using:

  \begin{itemize}
  \tightlist
  \item
    Analytic SFH models;
  \item
    Table-based SFH models;
  \item
    Particle-like data representations.
  \end{itemize}
\item
  Predicting observable quantities for a grid of models.
\end{itemize}

Full documentation is available
\href{https://population-synthesis-toolkit.readthedocs.io/en/latest/}{online}.

\hypertarget{acknowledgements}{%
\section{Acknowledgements}\label{acknowledgements}}

We acknowledge financial support from the Spanish State Research Agency
(AEI/10.13039/501100011033) through grant PID2019-107408GB-C42.

Daniel Jim\'enez-L\'opez was supported by Fondo Europeo de Desarrollo
Regional (MINCIN/AEI/10.13039/501100011033/FEDER, UE), through a
FPI-contract fellowship in the project PID2022-138560NB.

Our package relies on several widely used open-source Python libraries,
including \href{http://www.numpy.org}{NumPy} (Harris et al., 2020),
\href{https://www.matplotlib.org/}{Matplotlib} (Hunter, 2007) and
\href{http://www.astropy.org}{Astropy} (Astropy Collaboration et al.,
2022, 2018, 2013).

\hypertarget{references}{%
\section*{References}\label{references}}
\addcontentsline{toc}{section}{References}

\hypertarget{refs}{}
\begin{CSLReferences}{1}{0}
\leavevmode\hypertarget{ref-astropy:2022}{}%
Astropy Collaboration, Price-Whelan, A. M., Lim, P. L., Earl, N.,
Starkman, N., Bradley, L., Shupe, D. L., Patil, A. A., Corrales, L.,
Brasseur, C. E., N"othe, M., Donath, A., Tollerud, E., Morris, B. M.,
Ginsburg, A., Vaher, E., Weaver, B. A., Tocknell, J., Jamieson, W.,
\ldots{} Astropy Project Contributors. (2022). The {Astropy Project}:
Sustaining and growing a community-oriented open-source project and the
latest major release (v5.0) of the core package. \emph{The Astrophysical
Journal}, \emph{935}(2), 167.
\url{https://doi.org/10.3847/1538-4357/ac7c74}

\leavevmode\hypertarget{ref-astropy:2018}{}%
Astropy Collaboration, Price-Whelan, A. M., Sipőcz, B. M., Günther, H.
M., Lim, P. L., Crawford, S. M., Conseil, S., Shupe, D. L., Craig, M.
W., Dencheva, N., Ginsburg, A., Vand erPlas, J. T., Bradley, L. D.,
P\'erez-Su\'arez, D., de Val-Borro, M., Aldcroft, T. L., Cruz, K. L.,
Robitaille, T. P., Tollerud, E. J., \ldots{} Astropy Contributors.
(2018). The {Astropy Project}: Building an open-science project and
status of the v2.0 core package. \emph{The Astronomical Journal},
\emph{156}(3), 123. \url{https://doi.org/10.3847/1538-3881/aabc4f}

\leavevmode\hypertarget{ref-astropy:2013}{}%
Astropy Collaboration, Robitaille, T. P., Tollerud, E. J., Greenfield,
P., Droettboom, M., Bray, E., Aldcroft, T., Davis, M., Ginsburg, A.,
Price-Whelan, A. M., Kerzendorf, W. E., Conley, A., Crighton, N.,
Barbary, K., Muna, D., Ferguson, H., Grollier, F., Parikh, M. M., Nair,
P. H., \ldots{} Streicher, O. (2013). Astropy: A community {Python}
package for astronomy. \emph{Astronomy \& Astrophysics}, \emph{558},
A33. \url{https://doi.org/10.1051/0004-6361/201322068}

\leavevmode\hypertarget{ref-boquien+19}{}%
Boquien, M., Burgarella, D., Roehlly, Y., Buat, V., Ciesla, L., Corre,
D., Inoue, A. K., \& Salas, H. (2019). {CIGALE: a python Code
Investigating GALaxy Emission}. \emph{Astronomy \& Astrophysics},
\emph{622}, A103. \url{https://doi.org/10.1051/0004-6361/201834156}

\leavevmode\hypertarget{ref-bc+03}{}%
Bruzual, G., \& Charlot, S. (2003). {Stellar population synthesis at the
resolution of 2003}. \emph{Monthly Notices of the Royal Astronomical
Society}, \emph{344}(4), 1000--1028.
\url{https://doi.org/10.1046/j.1365-8711.2003.06897.x}

\leavevmode\hypertarget{ref-capellari+04}{}%
Cappellari, M., \& Emsellem, E. (2004). Parametric recovery of
line-of-sight velocity distributions from absorption-line spectra of
galaxies via penalized likelihood. \emph{Publications of the
Astronomical Society of the Pacific}, \emph{116}(816), 138--147.
\url{https://doi.org/10.1086/381875}

\leavevmode\hypertarget{ref-cid-fernandes+05}{}%
Cid Fernandes, R., Mateus, A., Sodr\'e, L., Stasi\'nska, G., \& Gomes, J. M.
(2005). {Semi-empirical analysis of Sloan Digital Sky Survey galaxies -
I. Spectral synthesis method}. \emph{Monthly Notices of the Royal
Astronomical Society}, \emph{358}(2), 363--378.
\url{https://doi.org/10.1111/j.1365-2966.2005.08752.x}

\leavevmode\hypertarget{ref-conroyux5cux26gunn10}{}%
Conroy, C., \& Gunn, J. E. (2010). The propagation of uncertainties in
stellar population synthesis modeling. {III}. Model calibration,
comparison, and evaluation. \emph{The Astrophysical Journal},
\emph{712}(2), 833--857.
\url{https://doi.org/10.1088/0004-637X/712/2/833}

\leavevmode\hypertarget{ref-conroy+09}{}%
Conroy, C., Gunn, J. E., \& White, M. (2009). The propagation of
uncertainties in stellar population synthesis modeling. {I}. The
relevance of uncertain aspects of stellar evolution and the initial mass
function to the derived physical properties of galaxies. \emph{The
Astrophysical Journal}, \emph{699}(1), 486--506.
\url{https://doi.org/10.1088/0004-637X/699/1/486}

\leavevmode\hypertarget{ref-pykoala_adass}{}%
Corcho-Caballero, P., Ascasibar, Y., L\'opez-S\'anchez, A. R.,
Gonz\'alez-Bol\'ivar, M., Lorente, N., Tocknell, J., Jim\'enez-Ibarra, F.,
Daluwathumullagamage, P. J., Quattropani, G., Owers, M., \&
Verdoes-Kleijn, G. A. (2025). \emph{The {PyKOALA} {Python} library: A
multi-instrument package for {IFS} data reduction}.
\url{https://arxiv.org/abs/2507.18347}

\leavevmode\hypertarget{ref-cc+25}{}%
Euclid Collaboration: Corcho-Caballero, P., Ascasibar, Y., Verdoes
Kleijn, G., Lovell, C. C., De Lucia, G., Cleland, C., Fontanot, F.,
Tortora, C., Koopmans, L. V. E., Moutard, T., Laigle, C., Nersesian, A.,
Shankar, F., Aghanim, N., Altieri, B., Amara, A., Andreon, S., Aussel,
H., Baccigalupi, C., \ldots{} Tenti, M. (2025). {Euclid Quick Data
Release (Q1). A probabilistic classification of quenched galaxies}.
\emph{arXiv e-Prints}, arXiv:2503.15315.
\url{https://doi.org/10.48550/arXiv.2503.15315}

\leavevmode\hypertarget{ref-harris2020array}{}%
Harris, C. R., Millman, K. J., Walt, S. J. van der, Gommers, R.,
Virtanen, P., Cournapeau, D., Wieser, E., Taylor, J., Berg, S., Smith,
N. J., Kern, R., Picus, M., Hoyer, S., Kerkwijk, M. H. van, Brett, M.,
Haldane, A., R\'io, J. F. del, Wiebe, M., Peterson, P., \ldots{} Oliphant,
T. E. (2020). Array programming with {NumPy}. \emph{Nature},
\emph{585}(7825), 357--362.
\url{https://doi.org/10.1038/s41586-020-2649-2}

\leavevmode\hypertarget{ref-hearin+23}{}%
Hearin, A. P., Chaves-Montero, J., Alarcon, A., Becker, M. R., \&
Benson, A. (2023). {DSPS: Differentiable stellar population synthesis}.
\emph{Monthly Notices of the Royal Astronomical Society}, \emph{521}(2),
1741--1756. \url{https://doi.org/10.1093/mnras/stad456}

\leavevmode\hypertarget{ref-hendriksux5cux26Izzard23}{}%
Hendriks, D., \& Izzard, R. (2023). {binary\_c-python: A Python-based
stellar population synthesis tool and interface to binary\_c}. \emph{The
Journal of Open Source Software}, \emph{8}(85), 4642.
\url{https://doi.org/10.21105/joss.04642}

\leavevmode\hypertarget{ref-hosek+20}{}%
Hosek, M. W., Jr., Lu, J. R., Lam, C. Y., Gautam, A. K., Lockhart, K.
E., Kim, D., \& Jia, S. (2020). {SPISEA}: A {Python}-based simple
stellar population synthesis code for star clusters. \emph{The
Astronomical Journal}, \emph{160}(3), 143.
\url{https://doi.org/10.3847/1538-3881/aba533}

\leavevmode\hypertarget{ref-hunter:2007}{}%
Hunter, J. D. (2007). Matplotlib: A {2D} graphics environment.
\emph{Computing in Science \& Engineering}, \emph{9}(3), 90--95.
\url{https://doi.org/10.1109/MCSE.2007.55}

\leavevmode\hypertarget{ref-jonhson+24}{}%
Johnson, B. D., Foreman-Mackey, D., Sick, J., Leja, J., Walmsley, M.,
Tollerud, E., Leung, H., Scott, S., \& Park, M. (2024).
\emph{{dfm/python-fsps: v0.4.7}} (Version v0.4.7) {[}Computer
software{]}. Zenodo. \url{https://doi.org/10.5281/zenodo.12447779}

\leavevmode\hypertarget{ref-johnson+21}{}%
Johnson, B. D., Leja, J., Conroy, C., \& Speagle, J. S. (2021). Stellar
population inference with {Prospector}. \emph{The Astrophysical Journal
Supplement Series}, \emph{254}(2), 22.
\url{https://doi.org/10.3847/1538-4365/abef67}

\leavevmode\hypertarget{ref-molla+09}{}%
Moll\'a, M., Garcı\'ia-Vargas, M. L., \& Bressan, A. (2009). {PopStar I}:
Evolutionary synthesis model description. \emph{Monthly Notices of the
Royal Astronomical Society}, \emph{398}(1), 451--470.
\url{https://doi.org/10.1111/j.1365-2966.2009.15160.x}

\leavevmode\hypertarget{ref-robotham+20}{}%
Robotham, A. S. G., Bellstedt, S., Lagos, C. del P., Thorne, J. E.,
Davies, L. J., Driver, S. P., \& Bravo, M. (2020). {ProSpect}:
Generating spectral energy distributions with complex star formation and
metallicity histories. \emph{Monthly Notices of the Royal Astronomical
Society}, \emph{495}(1), 905--931.
\url{https://doi.org/10.1093/mnras/staa1116}

\leavevmode\hypertarget{ref-rodrigo+24}{}%
Rodrigo, C., Cruz, P., Aguilar, J. F., Aller, A., Solano, E.,
G\'alvez-Ortiz, M. C., Jim\'enez-Esteban, F., Mas-Buitrago, P., Bayo, A.,
Cort\'es-Contreras, M., Murillo-Ojeda, R., Bonoli, S., Cenarro, J., Dupke,
R., L\'opez-Sanjuan, C., Mar\'in-Franch, A., de Oliveira, C. M., Moles, M.,
Taylor, K., \ldots{} Rami\'o, H. V. (2024). {Photometric segregation of
dwarf and giant FGK stars using the SVO Filter Profile Service and
photometric tools}. \emph{Astronomy \& Astrophysics}, \emph{689}, A93.
\url{https://doi.org/10.1051/0004-6361/202449998}

\leavevmode\hypertarget{ref-rodrigo+20}{}%
Rodrigo, C., \& Solano, E. (2020). {The SVO Filter Profile Service}.
\emph{XIV.0 Scientific Meeting (Virtual) of the Spanish Astronomical
Society}, 182.

\leavevmode\hypertarget{ref-rodrigo+12}{}%
Rodrigo, C., Solano, E., \& Bayo, A. (2012). \emph{{SVO Filter Profile
Service Version 1.0}} (p. 1015). IVOA Working Draft 15 October 2012.
\url{https://doi.org/10.5479/ADS/bib/2012ivoa.rept.1015R}

\leavevmode\hypertarget{ref-sanchez+16}{}%
S\'anchez, S. F., P\'erez, E., S\'anchez-Bl\'azquez, P., Gonz\'alez, J. J.,
Ros\'ales-Ortega, F. F., Cano-D\'iaz, M., L\'opez-Cob\'a, C., Marino, R. A., Gil
de Paz, A., Moll\'a, M., L\'opez-S\'anchez, A. R., Ascasibar, Y., \&
Barrera-Ballesteros, J. (2016). {Pipe3D}, a pipeline to analyze integral
field spectroscopy data: {I}. New fitting philosophy of {FIT3D}.
\emph{Revista Mexicana de Astronom\'ia y Astrof\'isica}, \emph{52}, 21--53.
\url{https://doi.org/10.48550/arXiv.1509.08552}

\leavevmode\hypertarget{ref-vazdekis+16}{}%
Vazdekis, A., Koleva, M., Ricciardelli, E., Röck, B., \& Falc\'on-Barroso,
J. (2016). {UV}-extended {E-MILES} stellar population models: Young
components in massive early-type galaxies. \emph{Monthly Notices of the
Royal Astronomical Society}, \emph{463}(4), 3409--3436.
\url{https://doi.org/10.1093/mnras/stw2231}

\leavevmode\hypertarget{ref-verro+22}{}%
Verro, K., Trager, S. C., Peletier, R. F., Lançon, A., Arentsen, A.,
Chen, Y.-P., Coelho, P. R. T., Dries, M., Falc\'on-Barroso, J., Gonneau,
A., Lyubenova, M., Martins, L., Prugniel, P., S\'anchez-Bl\'azquez, P., \&
Vazdekis, A. (2022). {Modelling simple stellar populations in the
near-ultraviolet to near-infrared with the X-shooter Spectral Library
(XSL)}. \emph{Astronomy \& Astrophysics}, \emph{661}, A50.
\url{https://doi.org/10.1051/0004-6361/202142387}

\leavevmode\hypertarget{ref-worthey+94}{}%
Worthey, G., Faber, S. M., Gonzalez, J. J., \& Burstein, D. (1994). Old
stellar populations. {V}. Absorption feature indices for the complete
{Lick/IDS} sample of stars. \emph{The Astrophysical Journal Supplement
Series}, \emph{94}, 687. \url{https://doi.org/10.1086/192087}

\leavevmode\hypertarget{ref-zuntz+15}{}%
Zuntz, J., Paterno, M., Jennings, E., Rudd, D., Manzotti, A., Dodelson,
S., Bridle, S., Sehrish, S., \& Kowalkowski, J. (2015). {CosmoSIS}:
Modular cosmological parameter estimation. \emph{Astronomy and
Computing}, \emph{12}, 45--59.
\url{https://doi.org/10.1016/j.ascom.2015.05.005}

\end{CSLReferences}

\end{document}
